\chapter{El teorema de Bolzano-Weierstrass}

\section{Introducción: una piedra angular del análisis}

El teorema de Bolzano-Weierstrass es uno de los resultados más célebres
del análisis real. Establece que toda sucesión acotada tiene una
subsucesión convergente. Más precisamente, en el marco que hemos
desarrollado:

\begin{quote}
  \itshape
  Toda sucesión acotada tiene una subsucesión de Cauchy.
\end{quote}

Este teorema es notable porque garantiza la existencia de comportamiento
convergente incluso en sucesiones que pueden parecer caóticas u
oscilatorias. Constituye el fundamento para:

\begin{itemize}
  \item Demostrar que las funciones continuas en intervalos cerrados
        alcanzan sus extremos.
  \item Establecer la compacidad de los subconjuntos cerrados y acotados
        de $\mathbb{R}^n$.
  \item Argumentos de compacidad secuencial a lo largo del análisis.
  \item Pruebas de existencia en optimización y ecuaciones diferenciales.
\end{itemize}

En esta capítulo demostraremos el Teorema de Bolzano-Weierstrass
combinando todas las herramientas desarrolladas: el teorema de la
subsucesión monótona del capítulo~13, el resultado de que las
sucesiones monótonas acotadas son de Cauchy del capítulo~12, y el
resultado correspondiente para sucesiones antótonas de las relaciones de
problemas 12 y~13.

\section{Nuevos teoremas}

\subsection{\texttt{abs\_le}: El complemento de \texttt{abs\_lt}}

Previamente hemos visto \texttt{abs\_lt}, que caracteriza cuándo
$|x| < y$. Ahora necesitamos su complemento para desigualdades no
estrictas.

\begin{theorem}[\texttt{abs\_le}]
  Para números reales o racionales $x, y$ con $y \geq 0$:
  \[
    |x| \leq y \iff -y \leq x \leq y.
  \]
\end{theorem}

Esta bicondicional nos permite separar las desigualdades de valor
absoluto en cotas simultáneas. En la dirección directa, $|x| \leq y$
implica tanto $x \leq y$ como $-x \leq y$, lo que es equivalente a
$-y \leq x$. En la dirección inversa, si $-y \leq x \leq y$, entonces
tanto $x$ como $-x$ están acotados por $y$, de modo que $|x| \leq y$.

\subsection{\texttt{IsCauchyOfAntitoneBdd}: El dual de nuestro resultado anterior}

En el capítulo~12 demostramos que las sucesiones monótonas
(no decrecientes) acotadas son de Cauchy. El resultado dual trata las
sucesiones no crecientes.

\begin{theorem}[\texttt{IsCauchyOfAntitoneBdd}]
  Si una sucesión $a : \mathbb{N} \to X$ (donde $X$ es $\mathbb{Q}$ o
  $\mathbb{R}$) es antótona (no creciente) y está acotada inferiormente
  por algún $M$, entonces $a$ es de Cauchy.
\end{theorem}

La técnica de demostración es elegante: se niega la sucesión para
convertir el problema antótono en uno monótono. Si $a$ es antótona y
acotada inferiormente por $M$, entonces $b = -a$ es monótona y acotada
superiormente por $-M$. Por nuestro teorema anterior, $b$ es de Cauchy.
Entonces $a = -b$ también es de Cauchy (las sucesiones de Cauchy son
cerradas bajo la negación).

Este fue el contenido de la relación de problemas~12, y lo usaremos como
una caja negra aquí.

\subsection{\texttt{AntitoneSubseq\_of\_UnBddPeaks}: La mitad faltante}

En el capítulo~13, demostramos que las sucesiones sin picos no
acotados tienen subsucesiones monótonas. El resultado complementario
trata el caso opuesto.

\begin{theorem}[\texttt{AntitoneSubseq\_of\_UnBddPeaks}]
  Si una sucesión $a : \mathbb{N} \to X$ tiene picos no acotados,
  entonces existe una subsucesión $\sigma : \mathbb{N} \to \mathbb{N}$
  tal que:
  \begin{itemize}
    \item $\sigma$ es estrictamente creciente (es un \texttt{Subseq}).
    \item $a \circ \sigma$ es antótona (no creciente).
  \end{itemize}
\end{theorem}

La intuición es sencilla: si hay picos arbitrariamente grandes, podemos
elegir una subsucesión formada únicamente por picos. Dado que cada pico
domina todos los valores futuros, y nuestra subsucesión consiste en
picos ordenados por sus índices, los valores de la sucesión a lo largo
de esta subsucesión deben ser no crecientes.

Este es el contenido de la relación de problemas~13, y lo usaremos como
una caja negra en nuestra demostración.

\section{El teorema principal: Bolzano-Weierstrass}

Ahora estamos listos para enunciar y demostrar el resultado principal.

\begin{theorem}[\texttt{BolzanoWeierstrass}]
  Si una sucesión $a : \mathbb{N} \to X$ (donde $X$ es $\mathbb{Q}$ o
  $\mathbb{R}$) es acotada, entonces existe una subsucesión
  $\sigma : \mathbb{N} \to \mathbb{N}$ tal que:
  \begin{itemize}
    \item $\sigma$ es estrictamente creciente (es un \texttt{Subseq}).
    \item $a \circ \sigma$ es de Cauchy.
  \end{itemize}
\end{theorem}

\subsection{Visión estratégica}

La demostración se apoya en la dicotomía establecida en la
Capítulo~13: toda sucesión tiene picos no acotados o no los tiene.
Trataremos cada caso por separado:

\medskip
\noindent\textbf{Caso 1: La sucesión tiene picos no acotados.}
\begin{enumerate}
  \item Por \texttt{AntitoneSubseq\_of\_UnBddPeaks}, se extrae una
        subsucesión antótona $\sigma$.
  \item La acotación de $a$ nos da cotas sobre $a \circ \sigma$.
  \item Por \texttt{IsCauchyOfAntitoneBdd}, se concluye que
        $a \circ \sigma$ es de Cauchy.
\end{enumerate}

\medskip
\noindent\textbf{Caso 2: La sucesión no tiene picos no acotados.}
\begin{enumerate}
  \item Por \texttt{MonotoneSubseq\_of\_BddPeaks} del capítulo~13,
        se extrae una subsucesión monótona $\sigma$.
  \item La acotación de $a$ nos da cotas sobre $a \circ \sigma$.
  \item Por \texttt{IsCauchyOfMonotoneBdd} del capítulo~12, se
        concluye que $a \circ \sigma$ es de Cauchy.
\end{enumerate}

La idea clave es que en ambos casos obtenemos una subsucesión monótona
(o antótona) acotada, y ambos tipos están garantizados como de Cauchy
por nuestros teoremas anteriores.

\subsection{Comprendiendo la acotación}

Cuando decimos que una sucesión $a$ es acotada (\texttt{SeqBdd} $a$),
queremos decir que existe algún $M > 0$ tal que $|a(n)| \leq M$ para
todo $n$.

Usando \texttt{abs\_le}, podemos descomponer esto en dos cotas
separadas:
\begin{align*}
  &\text{Cota superior: } a(n) \leq M \text{ para todo } n,\\
  &\text{Cota inferior: } {-M} \leq a(n) \text{ para todo } n.
\end{align*}

Estas cotas se preservan al pasar a subsucesiones: si $|a(n)| \leq M$
para todo $n$, entonces $|a(\sigma(n))| \leq M$ para todo $n$ y
cualquier subsucesión $\sigma$.

\subsection{Demostración en Lean}

\begin{minted}{lean4}
Statement BolzanoWeierstrass ( a : N → X )
    ( ha : SeqBdd a )
    : ∃ σ, Subseq σ ∧ IsCauchy ( a ◦ σ) := by
  choose M Mpos hM using ha
  have aBddAbove : ∀ n , a n ≤ M := by
    intro n ;
    specialize hM n ;
    rewrite [ abs_le ] at hM ;
    apply hM .2
  have aBddBelow : ∀ n , -M ≤ a n := by
    intro n ;
    specialize hM n ;
    rewrite [ abs_le ] at hM ;
    apply hM .1
  by_cases hPeaks : UnBddPeaks a
  · choose σ σsubseq σanti using
        AntitoneSubseq_of_UnBddPeaks a hPeaks
    use σ
    split_ands
    · exact σsubseq
    · apply IsCauchyOfAntitoneBdd ( a ◦ σ) ( - M )
      · intro n
        change -M ≤ a (σ n )
        apply aBddBelow
      · apply σanti
  · choose σ σsubseq σmono using
        MonotoneSubseq_of_BddPeaks a hPeaks
    use σ
    split_ands
    · apply σsubseq
    · apply IsCauchyOfMonotoneBdd ( a ◦ σ) M
      · intro n
        change a (σ n ) ≤ M
        apply aBddAbove
      · apply σmono
\end{minted}

\subsection{Demostración en lenguaje natural}

\begin{proof}
Sea $a : \mathbb{N} \to X$ una sucesión acotada. Por la definición de
\texttt{SeqBdd}, existe $M > 0$ tal que $|a(n)| \leq M$ para todo $n$.

\medskip
\noindent\textbf{Extracción de las cotas.}
Usando \texttt{abs\_le}, reescribimos $|a(n)| \leq M$ como la
conjunción
\[
  -M \leq a(n) \leq M,
\]
lo que nos da:
\begin{itemize}
  \item Cota superior: $a(n) \leq M$ para todo $n$.
  \item Cota inferior: $-M \leq a(n)$ para todo $n$.
\end{itemize}

\medskip
\noindent\textbf{Análisis por casos.}
Por el principio del tercio excluido, o bien \texttt{UnBddPeaks}$(a)$
se cumple o no. Trataremos cada caso por separado.

\medskip
\noindent\textit{Caso 1: $a$ tiene picos no acotados.}
Supongamos que \texttt{UnBddPeaks}$(a)$ se cumple.

Por el teorema \texttt{AntitoneSubseq\_of\_UnBddPeaks}, existe una
subsucesión $\sigma : \mathbb{N} \to \mathbb{N}$ tal que $\sigma$ es
estrictamente creciente y $a \circ \sigma$ es antótona.

La subsucesión $a \circ \sigma$ hereda la cota inferior de $a$: para
todo $n$, $-M \leq a(\sigma(n))$.

Por lo tanto, $a \circ \sigma$ es antótona y acotada inferiormente por
$-M$. Por el teorema \texttt{IsCauchyOfAntitoneBdd}, concluimos que
$a \circ \sigma$ es de Cauchy. Así, $\sigma$ satisface las dos
propiedades requeridas.

\medskip
\noindent\textit{Caso 2: $a$ no tiene picos no acotados.}
Supongamos que $\neg$\,\texttt{UnBddPeaks}$(a)$ se cumple.

Por el teorema \texttt{MonotoneSubseq\_of\_BddPeaks} de la
Capítulo~13, existe una subsucesión $\sigma : \mathbb{N} \to
\mathbb{N}$ tal que $\sigma$ es estrictamente creciente y $a \circ
\sigma$ es monótona.

La subsucesión $a \circ \sigma$ hereda la cota superior de $a$: para
todo $n$, $a(\sigma(n)) \leq M$.

Por lo tanto, $a \circ \sigma$ es monótona y acotada superiormente por
$M$. Por el teorema \texttt{IsCauchyOfMonotoneBdd} de la
Capítulo~12, concluimos que $a \circ \sigma$ es de Cauchy. Así,
$\sigma$ satisface las dos propiedades requeridas.

\medskip
\noindent\textbf{Conclusión.}
En ambos casos hemos construido una subsucesión $\sigma$ estrictamente
creciente tal que $a \circ \sigma$ es de Cauchy. Esto completa la
demostración.
\end{proof}

\section{La arquitectura de la demostración}

Demos un paso atrás y apreciemos la estructura elegante de este
argumento. El Teorema de Bolzano-Weierstrass es la cúpula de un
edificio cuidadosamente construido.

\subsection{La cimentación: sucesiones monótonas}

En la base tenemos el hecho fundamental de que las sucesiones monótonas
acotadas son de Cauchy. Esto se demostró en el capítulo~12 mediante
un ingenioso argumento de contradicción: si una sucesión monótona
acotada no fuera de Cauchy, podríamos construir una sucesión infinita
de brechas que eventualmente forzarían a la sucesión a superar su cota.

El resultado dual para sucesiones antótonas se sigue mediante un
sencillo argumento de negación.

\subsection{La dicotomía: picos y valles}

El siguiente nivel introduce el concepto de picos y establece una
dicotomía: toda sucesión tiene picos no acotados o no los tiene.

\begin{itemize}
  \item \textbf{Picos no acotados:} La sucesión regresa continuamente a
        nuevos máximos locales de los que nunca se recupera.
  \item \textbf{Picos acotados:} A partir de cierto punto, la sucesión
        no tiene más picos; debe empezar a subir eventualmente.
\end{itemize}

\subsection{La extracción: de la dicotomía a la monotonía}

El Capítulo~13 mostró que las sucesiones sin picos no acotados
tienen subsucesiones monótonas. La demostración usó la hermosa
construcción de la órbita: iterando una función cuidadosamente
elegida, se extrajo una subsucesión que automáticamente evitaba todos
los picos y, por tanto, se veía forzada a ser monótona.

La relación de problemas~13 estableció el resultado dual: las sucesiones
con picos no acotados tienen subsucesiones antótonas. La demostración es
más sencilla aquí: basta elegir los propios picos como la subsucesión.

\subsection{La síntesis: Bolzano-Weierstrass}

Finalmente, combinamos todos estos resultados. Toda sucesión acotada cae
en uno de los dos casos de nuestra dicotomía. En ambos casos podemos
extraer una subsucesión monótona (o antótona), y dicha subsucesión
permanece acotada. Por nuestros resultados fundamentales, toda sucesión
monótona acotada es de Cauchy.

El teorema no surge de un único argumento ingenioso, sino del ensamblaje
cuidadoso de múltiples resultados, cada uno construido sobre el anterior.

\section{Por qué esto importa}

El Teorema de Bolzano-Weierstrass es fundamental para el análisis por
varias razones.

\subsection{Compacidad}

En topología, se dice que un espacio es \emph{secuencialmente compacto}
si toda sucesión tiene una subsucesión convergente. El Teorema de
Bolzano-Weierstrass muestra que los subconjuntos cerrados y acotados de
$\mathbb{R}$ son secuencialmente compactos. Este es un paso clave para
demostrar el Teorema de Heine-Borel, que caracteriza los subconjuntos
compactos del espacio euclídeo.

\subsection{Pruebas de existencia}

Muchos resultados de existencia en análisis dependen del Teorema de
Bolzano-Weierstrass. Por ejemplo:

\begin{itemize}
  \item \textbf{Teorema del Valor Extremo:} Una función continua en un
        intervalo cerrado alcanza su máximo y su mínimo.
  \item \textbf{Teorema de Bolzano:} Una función continua que cambia de
        signo debe tener un cero.
  \item \textbf{Optimización:} Las funciones continuas en conjuntos
        compactos alcanzan sus valores óptimos.
\end{itemize}

\subsection{Aproximación}

En el análisis numérico y las matemáticas aplicadas, el Teorema de
Bolzano-Weierstrass garantiza que las sucesiones acotadas de
aproximaciones tienen subsucesiones convergentes, proporcionando una
base para los métodos iterativos.

\subsection{Análisis funcional}

El teorema se extiende a espacios de dimensión infinita (con las
modificaciones apropiadas), constituyendo la base para comprender la
convergencia débil y débil$^*$ en espacios de Banach.

\section{Contexto histórico}

El teorema lleva el nombre de Bernard Bolzano (1781--1848) y Karl
Weierstrass (1815--1897), dos pioneros del análisis riguroso.

El trabajo de Bolzano en los primeros años del siglo~XIX fue
notablemente moderno, anticipando muchas ideas que solo se generalizarían
más tarde. Comprendió la importancia de demostrar rigurosamente los
resultados de existencia y reconoció que la intuición sobre continuidad
y límites necesitaba una formalización cuidadosa.

Weierstrass, trabajando varias décadas después, suele recibir el crédito
de haber puesto el análisis sobre una base rigurosa. Insistió en las
definiciones $\varepsilon$-$\delta$ y construyó famosos contraejemplos
(como una función continua en ningún punto diferenciable) que mostraron
las limitaciones de la intuición geométrica.

El Teorema de Bolzano-Weierstrass representa su perspicacia compartida:
acotación más elección infinita fuerza acumulación. Si se tienen
infinitos valores confinados en una región acotada, deben agruparse en
algún lugar.

\section{Extensiones y variantes}

El Teorema de Bolzano-Weierstrass tiene muchas generalizaciones.

\subsection{Dimensiones superiores}

En $\mathbb{R}^n$, toda sucesión acotada tiene una subsucesión
convergente. La demostración sigue el mismo patrón: se extraen
subsucesiones monótonas en cada coordenada, usando un argumento diagonal
para garantizar que converjan simultáneamente.

\subsection{Espacios métricos}

En espacios métricos completos, un conjunto $K$ es compacto si y solo si
toda sucesión en $K$ tiene una subsucesión convergente (cuyo límite está
en $K$). Esta caracterización de la compacidad es una de las más útiles
en el análisis.

\subsection{Convergencia débil}

En espacios de Banach de dimensión infinita, las sucesiones acotadas
pueden no tener subsucesiones norm-convergentes (la bola unitaria podría
no ser compacta). Sin embargo, por el Teorema de Banach-Alaoglu, las
sucesiones acotadas en el espacio dual tienen subsucesiones
débil$^*$-convergentes.

\subsection{Espacios topológicos}

El marco más general son los espacios topológicos arbitrarios, donde la
compacidad se puede definir sin referencia a sucesiones. La relación
entre compacidad secuencial y compacidad depende de la topología (son
equivalentes para espacios métricos).

\section{Técnicas de demostración: Lo que aprendimos}

Esta demostración exhibe varias técnicas matemáticas importantes.

\subsection{Análisis por casos con \texttt{by\_cases}}

La táctica \texttt{by\_cases} nos permite dividir una demostración según
si una proposición se cumple o no. Esto es la formalización de la
demostración por casos, permitiéndonos aprovechar el principio del
tercio excluido.

\subsection{Descomposición de valores absolutos}

El teorema \texttt{abs\_le} muestra cómo trabajar con desigualdades de
valor absoluto separándolas en cotas unilaterales simultáneas. Esta
técnica es omnipresente en el análisis.

\subsection{Argumentos de dualidad}

La demostración de \texttt{IsCauchyOfAntitoneBdd} ejemplifica una
poderosa técnica: convertir un problema sobre sucesiones antótonas en
uno sobre sucesiones monótonas mediante la negación. Esto reduce el
número de teoremas que necesitamos demostrar desde cero.

\subsection{Aprovechamiento de resultados previos}

La demostración de Bolzano-Weierstrass es breve precisamente porque
hemos construido la maquinaria necesaria en capítulos anteriores.
Esto demuestra la importancia de:

\begin{itemize}
  \item Descomponer demostraciones complejas en lemas.
  \item Identificar y demostrar resultados generales útiles.
  \item Estructurar una teoría para que los teoremas principales sean
        consecuencias naturales.
\end{itemize}

\section{Mirando hacia adelante}

Con el teorema de Bolzano-Weierstrass en mano, estamos ahora equipados
para demostrar muchos resultados clásicos del análisis:

\begin{itemize}
  \item \textbf{Teorema del valor extremo:} las funciones continuas en
        intervalos cerrados son acotadas y alcanzan sus cotas.
  \item \textbf{Teorema del valor intermedio:} las funciones continuas
        en intervalos toman todos los valores intermedios.
  \item \textbf{Teorema de Heine-Borel:} caracterización de los
        subconjuntos compactos del espacio euclídeo.
  \item Resultados de continuidad uniforme en conjuntos compactos.
  \item Teoremas de punto fijo.
\end{itemize}

Cada uno de estos resultados se apoya en la perspicacia fundamental
capturada por Bolzano-Weierstrass: en el reino de lo acotado, siempre
hay convergencia escondida en algún lugar, esperando ser extraída.

%%% Local Variables:
%%% mode: LaTeX
%%% TeX-master: "Analisis_en_Lean"
%%% End:
